\section{Bound states in one dimension (10/14)}

\referto{\tong~2.3, 7.5}

A bound-state wave function is square-normalizable:
\begin{equation}
 \int_{-\infty}^{\infty}|\psi(x)|^2\dif x<\infty.
\end{equation}
Consequently,
\begin{equation}
 \lim_{|x|\to\infty}\psi(x)=0.
\end{equation}
For a potential with a well-defined asymptotic value, a bound-state energy
lies below the asymptotic continuum threshold:
\begin{equation}
 E<\lim_{|x|\to\infty}V(x).
\end{equation}
Unlike scattering states characterized by a continuous wavevector, bound states have a discrete spectrum. 

\subsection{The Infinite Potential Well}

\subsubsection{Potential and boundary conditions}
Consider a box of width $L$ centered at the origin:
\begin{equation}
 V(x)=
 \begin{cases}
 0, & |x|<L/2,\\
 +\infty, & |x|>L/2.
 \end{cases}
\end{equation}
The wave function vanishes outside the box and satisfies
\begin{equation}
 \psi(-L/2)=\psi(L/2)=0.
\end{equation}
Inside the box, a general energy eigenfunction is
\begin{equation}
 \psi(x)=A e^{\imth px/\hbar}+B e^{-\imth px/\hbar}.
\end{equation}
The two boundary conditions give
\begin{align}
 A e^{\imth pL/(2\hbar)}+B e^{-\imth pL/(2\hbar)}&=0,\\
 A e^{-\imth pL/(2\hbar)}+B e^{\imth pL/(2\hbar)}&=0.
\end{align}
A nontrivial solution requires
\begin{equation}
 e^{2\imth pL/\hbar}=1,
\end{equation}
so the allowed momenta are quantized.

\subsubsection{Eigenfunctions and energies}
It is convenient to write the normalized eigenfunctions as
\begin{equation}
 \boxed{
 \psi_n(x)=\sqrt{\frac{2}{L}}
 \sin\!\left[\frac{n\pi}{L}\left(x+\frac{L}{2}\right)\right],
 \qquad n=1,2,3,\ldots,}
\end{equation}
for $|x|<L/2$, with $\psi_n(x)=0$ outside the well. The energies are
\begin{equation}
 \boxed{E_n=\frac{\pi^2\hbar^2}{2mL^2}n^2,
 \qquad n=1,2,3,\ldots.}
\end{equation}
The ground state has no interior node, the first excited state has one node,
and the state with quantum number $n$ has $n-1$ interior nodes.

\subsection{General Theorems for One-Dimensional Bound States}

\subsubsection{Nondegeneracy}
The bound-state spectrum of a one-dimensional Schr\"odinger Hamiltonian is
nondegenerate. Suppose $\psi_1$ and $\psi_2$ are square-normalizable
solutions with the same energy $E$:
\begin{align}
 -\frac{\hbar^2}{2m}\psi_1''+V\psi_1&=E\psi_1,\\
 -\frac{\hbar^2}{2m}\psi_2''+V\psi_2&=E\psi_2.
\end{align}
Multiply the first equation by $\psi_2$, the second by $\psi_1$, and
subtract. This gives
\begin{equation}
 \frac{\dif{}}{\dif x}\left(\psi_2\psi_1'-\psi_1\psi_2'\right)=0.
\end{equation}
Thus the Wronskian is constant. Bound-state boundary conditions make this
constant zero, so
\begin{equation}
 \psi_2\psi_1'=\psi_1\psi_2'.
\end{equation}
Where the wave functions are nonzero,
\begin{equation}
 \frac{\psi_1'}{\psi_1}=\frac{\psi_2'}{\psi_2},
\end{equation}
and hence
\begin{equation}
 \psi_1(x)=C\psi_2(x).
\end{equation}
The two solutions therefore represent the same physical state.

\subsubsection{Node theorems}
For one-dimensional bound states:
\begin{enumerate}[label=\arabic*.]
 \item The ground-state wave function has no node at finite $x$.
 \item The $m$th excited state has exactly $m$ nodes.
\end{enumerate}
Equivalently, if the eigenstates are indexed by $n=1,2,3,\ldots$, then the
$n$th state has $n-1$ nodes.

\subsection{The Finite Square Well}


Consider the symmetric finite well
\begin{equation}
 V(x)=
 \begin{cases}
 -V_0, & |x|<a,\\
 0, & |x|>a,
 \end{cases}
 \qquad V_0>0.
\end{equation}

\subsubsection{Scattering states with $E>0$}

For scattering states with $E>0$, define
\begin{equation}
 k=\frac{\sqrt{2mE}}{\hbar},
 \qquad
 k'=\frac{\sqrt{2m(E+V_0)}}{\hbar}.
\end{equation}
A left-incident scattering state has the form
\begin{equation}
 \psi(x)=
 \begin{cases}
 e^{\imth kx}+r e^{-\imth kx}, & x<-a,\\
 A e^{\imth k'x}+B e^{-\imth k'x}, & |x|<a,\\
 t e^{\imth kx}, & x>a.
 \end{cases}
\end{equation}
Continuity of $\psi$ and $\psi'$ at $x=-a$ gives
\begin{align}
 e^{-\imth ka}+r e^{\imth ka}
 &=A e^{-\imth k'a}+B e^{\imth k'a},\\
 k\left(e^{-\imth ka}-r e^{\imth ka}\right)
 &=k'\left(A e^{-\imth k'a}-B e^{\imth k'a}\right).
\end{align}
At $x=a$,
\begin{align}
 A e^{\imth k'a}+B e^{-\imth k'a}&=t e^{\imth ka},\\
 k'\left(A e^{\imth k'a}-B e^{-\imth k'a}\right)&=kt e^{\imth ka}.
\end{align}
These matching equations determine the reflection and transmission
amplitudes.

Solving the two matching equations at \(x=\pm a\) for \(A\) and \(B\), we find \begin{align}\label{1d finite well:mathcing A} A &= \frac{e^{\imth k'a}}{2} \left[ \left(1+\frac{k}{k'}\right)e^{-\imth ka} + r\left(1-\frac{k}{k'}\right)e^{\imth ka} \right]= \frac{e^{-\imth k'a}}{2}\, t e^{\imth ka} \left(1+\frac{k}{k'}\right), \\ B &= \frac{e^{-\imth k'a}}{2} \left[ \left(1-\frac{k}{k'}\right)e^{-\imth ka} + r\left(1+\frac{k}{k'}\right)e^{\imth ka} \right]= \frac{e^{\imth k'a}}{2}\, t e^{\imth ka} \left(1-\frac{k}{k'}\right). \label{1d finite well:mathcing B}\end{align}


\subsubsection{Bound states with $-V_0<E<0$}
For bound states,
\begin{equation}
 -V_0<E<0.
\end{equation}
Define
\begin{equation}
 \kappa=-\imth k=\frac{\sqrt{2m|E|}}{\hbar},
 \qquad
 k'=\frac{\sqrt{2m(E+V_0)}}{\hbar}.
\end{equation}
These obey
\begin{equation}
 \kappa^2+(k')^2=\frac{2mV_0}{\hbar^2}.
\end{equation}
Normalizability requires exponential decay outside the well:
\begin{equation}
 \psi(x)=
 \begin{cases}
 r e^{\kappa x}, & x<-a,\\
 A e^{\imth k'x}+B e^{-\imth k'x}, & |x|<a,\\
 t e^{-\kappa x}, & x>a.
 \end{cases}
\end{equation}
Note that the matching conditions for bound states are exactly the same as for scattering states, except for replacing $k=\imth\kappa$. Comparing the form of wavefunctions, we see that in bound states
\begin{align}
    |r|,|t|\rightarrow\infty
\end{align}
and Eqs. (\ref{1d finite well:mathcing A})-(\ref{1d finite well:mathcing B}) becomes
\begin{equation} 
\frac{r}{t} = e^{2\imth k'a}\frac{k'-k}{k'+k} = e^{-2\imth k'a}\frac{k'+k}{k'-k} = \pm 1. \end{equation}
for bound states. 


\subsubsection{Inversion symmetry}
For an even potential,
\begin{equation}
 V(x)=V(-x),
\end{equation}
define the parity operator
\begin{equation}
 \hat P=\int_{-\infty}^{\infty}\dif x\,\dket{-x}\dbra{x}.
\end{equation}
It satisfies
\begin{equation}
 \hat P^2=\id,
 \qquad
 \hat P\hat x\hat P^{-1}=-\hat x,
 \qquad
 \hat P\hat p\hat P^{-1}=-\hat p.
\end{equation}
For
\begin{equation}
 \hat H=\frac{\hat p^{\,2}}{2m}+V(\hat x),
\end{equation}
one has
\begin{equation}
 \comm{\hat P}{\hat H}=0.
\end{equation}
Because the one-dimensional bound spectrum is nondegenerate, each bound
energy eigenstate is also a parity eigenstate:
\begin{equation}
 \hat P\dket{\psi}=\frac{r}{t}\dket\psi=\pm\dket{\psi}.
\end{equation}
Thus the wave function is either even or odd under the parity operator, depending on the sign of
\begin{align}\label{1d finite well:parity}
    \frac rt=e^{-2\imth k'a}\frac{k'+\imth\kappa}{k'-\imth\kappa}=\pm1,~~~~\kappa=\sqrt{2mV_0/\hbar^2-(k')^2}
\end{align}

\subsubsection{Even-parity bound states}
For an even state, choose
\begin{equation}
 \psi(x)=
 \begin{cases}
 C e^{\kappa x}, & x<-a,\\
 D\cos(k'x), & |x|<a,\\
 C e^{-\kappa x}, & x>a.
 \end{cases}
\end{equation}
Matching $\psi$ and $\psi'$ at $x=a$ gives
\begin{equation}
 \boxed{k'\tan(k'a)=\kappa.}
\end{equation}
Using $\kappa^2+(k')^2=2mV_0/\hbar^2$,
\begin{equation}
 \boxed{k'\tan(k'a)
 =\sqrt{\frac{2mV_0}{\hbar^2}-(k')^2}.}
\end{equation}

\subsubsection{Odd-parity bound states}
For an odd state, choose
\begin{equation}
 \psi(x)=
 \begin{cases}
 -C e^{\kappa x}, & x<-a,\\
 D\sin(k'x), & |x|<a,\\
 C e^{-\kappa x}, & x>a.
 \end{cases}
\end{equation}
Matching at $x=a$ gives
\begin{equation}
 \boxed{-k'\cot(k'a)=\kappa,}
\end{equation}
or
\begin{equation}
 \boxed{-k'\cot(k'a)
 =\sqrt{\frac{2mV_0}{\hbar^2}-(k')^2}.}
\end{equation}
Equivalently,
\begin{equation}
 \tan\!\left(k'a+\frac{\pi}{2}\right)
 =\frac{\kappa}{k'}.
\end{equation}

\subsubsection{Graphical solution}
The allowed bound-state energies are obtained by solving Eq. (\ref{1d finite well:parity}), which can be simplified as
\begin{align}
    k'\tan(k'a)=\sqrt{\frac{r}{t}}\cdot\sqrt{\frac{2mV_0}{\hbar^2}-(k')^2}
\end{align}
Graphically, as shown in Fig. \ref{fig:finite-well-graphical-solution}, the allowed values of $k'$ are the intersections of the red curve
\begin{equation}
 \kappa(k')=\sqrt{\frac{2mV_0}{\hbar^2}-(k')^2}
\end{equation}
with the green even-parity branches $k'\tan(k'a)$ and the blue odd-parity
branches $-k'\cot(k'a)$. Intersections with $k'\tan(k'a)$ give even-parity states, while intersections
with $-k'\cot(k'a)$ give odd-parity states. The vertical dashed lines indicate the poles at
$k'a=\pi/2,\pi,\ldots$, and the red curve meets the horizontal axis at
\begin{equation}
 k'a=a\sqrt{\frac{2mV_0}{\hbar^2}}.
\end{equation}

\begin{figure}[htbp]
 \centering
 \includegraphics[width=\textwidth]{1d_finite_well_bound_states.png}
 \caption{Graphical solution of the even- and odd-parity finite-well
 bound-state equations. Intersections with the red curve determine the
 allowed bound-state energies.}
 \label{fig:finite-well-graphical-solution}
\end{figure}


Once an allowed $k'$ is found,
\begin{equation}
 E=-V_0+\frac{\hbar^2(k')^2}{2m}.
\end{equation}
Only finitely many intersections occur for a well of finite depth and width,
so a finite square well supports finitely many bound states.

Finally, note that in the limit of 
\begin{align}
    a\rightarrow0^+,~~~V_0\rightarrow+\infty,~~~2aV_0\rightarrow|D_0|
\end{align}
the square potential well discussed here reduces to the delta function potential studied in the last section. 

\subsection{Summary}
%\addcontentsline{toc}{chapter}{Summary}
\begin{itemize}
 \item Bound states are square-normalizable and decay at spatial infinity.
 \item One-dimensional bound-state energy levels are nondegenerate.
 \item The ground state has no nodes, and successive excited states acquire
       one additional node each.
 \item For an even potential, nondegenerate bound states have definite parity.
 \item Finite-well energies follow from separate even- and odd-parity
       transcendental equations.
\end{itemize}

